
\section{Discussion}

\subsection{Key Takeaways}

\subsubsection{Takeaway 1: Cosformer is a Strong Baseline}

Of the linear attention mechanisms we tested, Cosformer is among the best 
at preserving accuracy. The cos/sin reweighting genuinely recovers locality.

\subsubsection{Takeaway 2: Encoder vs Decoder Paths Diverge}

The encoder path is clean and fast. The decoder path adds overhead from 
the chunked scan. On encoders (DistilBERT), Cosformer is closer to baseline.

\subsubsection{Takeaway 3: RoBERTa Needs Special Handling}

The nested encoder structure requires careful handling. Our default 
hyperparameters are insufficient; lower LR or excluding cos/sin reweighting 
helps.

\subsection{When to Use Cosformer}

\subsubsection{Recommended}

\begin{itemize}
\item \textbf{Long-context inference on flat encoders} (e.g., DistilBERT).
\item \textbf{Research baseline} for linear attention comparisons.
\item \textbf{When memory is the bottleneck} -- $\mathcal{O}(L)$ state.

\end{itemize}
\subsubsection{Not Recommended}

\begin{itemize}
\item \textbf{RoBERTa or other nested encoders} without hyperparameter tuning.
\item \textbf{Latency-critical inference} -- chunked scan adds overhead.
\item \textbf{Tasks requiring sharp attention} -- Cosformer smooths attention.

\end{itemize}
\subsection{Decision Flowchart}

\begin{lstlisting}
Start
  |
   
Is the model a flat encoder? -- Yes -- Use Cosformer
  |
  No
   
Is the model a decoder? -- Yes -- Consider but expect 3-5$\times$ slowdown
  |
  No
   
Is the model RoBERTa? -- Yes -- Tune LR first
  |
  No
   
Consider alternatives (FlashAttention, A3)
\end{lstlisting}

\subsection{Cost-Benefit Analysis}

For DistilBERT on RTX 4070:

\begin{table}[H]
\centering
\small
\begin{tabular}{llll}
\toprule
\textbf{Metric} & \textbf{Baseline} & \textbf{Cosformer} & \textbf{Change} \\
\midrule
Accuracy & 90.71\% & 77.75\% & -12.96\% \\
Params & 66.9M & 70791394M & +5.73\% \\
Time & 1.9 s & 3.5 s & 1.84$\times$ \\
VRAM & 183 MB & 335 MB MB & 1.31$\times$ \\
\bottomrule
\end{tabular}
\end{table}

\textbf{Verdict.} works on flat encoder; problematic on RoBERTa; slow on decoder

\subsection{Comparison with Alternatives}

\begin{table}[H]
\centering
\small
\begin{tabular}{llll}
\toprule
\textbf{Alternative} & \textbf{Accuracy} & \textbf{Speed} & \textbf{Params} \\
\midrule
FlashAttention-2 & ~90\% & 1.1$\times$ & 0 \\
A3 (SVD) & 88.5\% & 0.8$\times$ & \textbf{-60\%} \\
NHA & 90.25\% & 0.6$\times$ & +0.88\% \\
\textbf{Cosformer} & \textbf{77.75\%} & \textbf{1.84$\times$} & \textbf{+5.73\%} \\
\bottomrule
\end{tabular}
\end{table}

\subsection{Open Questions}

\begin{enumerate}
\item \textbf{Why does RoBERTa fail?} Our analysis is incomplete.
\item \textbf{Would larger LoRA rank help RoBERTa?} Untested.
\item \textbf{Would removing cos/sin for RoBERTa help?} Partially tested.
\item \textbf{How does Cosformer scale to L > 512?} Untested.

\end{enumerate}
\subsection{Broader Impact}

\subsubsection{Positive}

\begin{itemize}
\item Linear attention is more memory-efficient.
\item Cosformer is a well-tested baseline.

\end{itemize}
\subsubsection{Negative}

\begin{itemize}
\item Our implementation is slower than baseline on short sequences.
\item RoBERTa results are negative.

\end{itemize}
\subsection{Deployment Recommendations}

\textbf{Test on your specific model before deploying.}
\begin{itemize}
\item For flat encoders: likely fine.
\item For nested encoders: tune first.
\item For decoders: expect 3-5$\times$ slowdown.
\end{itemize}
\begin{lstlisting}

---
\end{lstlisting}